% Options for packages loaded elsewhere
\PassOptionsToPackage{unicode}{hyperref}
\PassOptionsToPackage{hyphens}{url}
\documentclass[
]{article}
\usepackage{xcolor}
\usepackage[margin=1in]{geometry}
\usepackage{amsmath,amssymb}
\setcounter{secnumdepth}{-\maxdimen} % remove section numbering
\usepackage{iftex}
\ifPDFTeX
  \usepackage[T1]{fontenc}
  \usepackage[utf8]{inputenc}
  \usepackage{textcomp} % provide euro and other symbols
\else % if luatex or xetex
  \usepackage{unicode-math} % this also loads fontspec
  \defaultfontfeatures{Scale=MatchLowercase}
  \defaultfontfeatures[\rmfamily]{Ligatures=TeX,Scale=1}
\fi
\usepackage{lmodern}
\ifPDFTeX\else
  % xetex/luatex font selection
\fi
% Use upquote if available, for straight quotes in verbatim environments
\IfFileExists{upquote.sty}{\usepackage{upquote}}{}
\IfFileExists{microtype.sty}{% use microtype if available
  \usepackage[]{microtype}
  \UseMicrotypeSet[protrusion]{basicmath} % disable protrusion for tt fonts
}{}
\makeatletter
\@ifundefined{KOMAClassName}{% if non-KOMA class
  \IfFileExists{parskip.sty}{%
    \usepackage{parskip}
  }{% else
    \setlength{\parindent}{0pt}
    \setlength{\parskip}{6pt plus 2pt minus 1pt}}
}{% if KOMA class
  \KOMAoptions{parskip=half}}
\makeatother
\usepackage{color}
\usepackage{fancyvrb}
\newcommand{\VerbBar}{|}
\newcommand{\VERB}{\Verb[commandchars=\\\{\}]}
\DefineVerbatimEnvironment{Highlighting}{Verbatim}{commandchars=\\\{\}}
% Add ',fontsize=\small' for more characters per line
\usepackage{framed}
\definecolor{shadecolor}{RGB}{248,248,248}
\newenvironment{Shaded}{\begin{snugshade}}{\end{snugshade}}
\newcommand{\AlertTok}[1]{\textcolor[rgb]{0.94,0.16,0.16}{#1}}
\newcommand{\AnnotationTok}[1]{\textcolor[rgb]{0.56,0.35,0.01}{\textbf{\textit{#1}}}}
\newcommand{\AttributeTok}[1]{\textcolor[rgb]{0.13,0.29,0.53}{#1}}
\newcommand{\BaseNTok}[1]{\textcolor[rgb]{0.00,0.00,0.81}{#1}}
\newcommand{\BuiltInTok}[1]{#1}
\newcommand{\CharTok}[1]{\textcolor[rgb]{0.31,0.60,0.02}{#1}}
\newcommand{\CommentTok}[1]{\textcolor[rgb]{0.56,0.35,0.01}{\textit{#1}}}
\newcommand{\CommentVarTok}[1]{\textcolor[rgb]{0.56,0.35,0.01}{\textbf{\textit{#1}}}}
\newcommand{\ConstantTok}[1]{\textcolor[rgb]{0.56,0.35,0.01}{#1}}
\newcommand{\ControlFlowTok}[1]{\textcolor[rgb]{0.13,0.29,0.53}{\textbf{#1}}}
\newcommand{\DataTypeTok}[1]{\textcolor[rgb]{0.13,0.29,0.53}{#1}}
\newcommand{\DecValTok}[1]{\textcolor[rgb]{0.00,0.00,0.81}{#1}}
\newcommand{\DocumentationTok}[1]{\textcolor[rgb]{0.56,0.35,0.01}{\textbf{\textit{#1}}}}
\newcommand{\ErrorTok}[1]{\textcolor[rgb]{0.64,0.00,0.00}{\textbf{#1}}}
\newcommand{\ExtensionTok}[1]{#1}
\newcommand{\FloatTok}[1]{\textcolor[rgb]{0.00,0.00,0.81}{#1}}
\newcommand{\FunctionTok}[1]{\textcolor[rgb]{0.13,0.29,0.53}{\textbf{#1}}}
\newcommand{\ImportTok}[1]{#1}
\newcommand{\InformationTok}[1]{\textcolor[rgb]{0.56,0.35,0.01}{\textbf{\textit{#1}}}}
\newcommand{\KeywordTok}[1]{\textcolor[rgb]{0.13,0.29,0.53}{\textbf{#1}}}
\newcommand{\NormalTok}[1]{#1}
\newcommand{\OperatorTok}[1]{\textcolor[rgb]{0.81,0.36,0.00}{\textbf{#1}}}
\newcommand{\OtherTok}[1]{\textcolor[rgb]{0.56,0.35,0.01}{#1}}
\newcommand{\PreprocessorTok}[1]{\textcolor[rgb]{0.56,0.35,0.01}{\textit{#1}}}
\newcommand{\RegionMarkerTok}[1]{#1}
\newcommand{\SpecialCharTok}[1]{\textcolor[rgb]{0.81,0.36,0.00}{\textbf{#1}}}
\newcommand{\SpecialStringTok}[1]{\textcolor[rgb]{0.31,0.60,0.02}{#1}}
\newcommand{\StringTok}[1]{\textcolor[rgb]{0.31,0.60,0.02}{#1}}
\newcommand{\VariableTok}[1]{\textcolor[rgb]{0.00,0.00,0.00}{#1}}
\newcommand{\VerbatimStringTok}[1]{\textcolor[rgb]{0.31,0.60,0.02}{#1}}
\newcommand{\WarningTok}[1]{\textcolor[rgb]{0.56,0.35,0.01}{\textbf{\textit{#1}}}}
\usepackage{graphicx}
\makeatletter
\newsavebox\pandoc@box
\newcommand*\pandocbounded[1]{% scales image to fit in text height/width
  \sbox\pandoc@box{#1}%
  \Gscale@div\@tempa{\textheight}{\dimexpr\ht\pandoc@box+\dp\pandoc@box\relax}%
  \Gscale@div\@tempb{\linewidth}{\wd\pandoc@box}%
  \ifdim\@tempb\p@<\@tempa\p@\let\@tempa\@tempb\fi% select the smaller of both
  \ifdim\@tempa\p@<\p@\scalebox{\@tempa}{\usebox\pandoc@box}%
  \else\usebox{\pandoc@box}%
  \fi%
}
% Set default figure placement to htbp
\def\fps@figure{htbp}
\makeatother
\setlength{\emergencystretch}{3em} % prevent overfull lines
\providecommand{\tightlist}{%
  \setlength{\itemsep}{0pt}\setlength{\parskip}{0pt}}
\usepackage{bookmark}
\IfFileExists{xurl.sty}{\usepackage{xurl}}{} % add URL line breaks if available
\urlstyle{same}
\hypersetup{
  hidelinks,
  pdfcreator={LaTeX via pandoc}}

\author{}
\date{\vspace{-2.5em}}

\begin{document}

\section{Replicator Dynamics}\label{sec-replicator-dynamics}

\begin{quote}
An introduction to the replicator equation --- the central dynamical
system of evolutionary game theory --- with ODE solutions in R for
Rock-Paper-Scissors and Hawk-Dove games, including simplex
visualization.
\end{quote}

\subsection*{Learning objectives}\label{learning-objectives}
\addcontentsline{toc}{subsection}{Learning objectives}

\begin{itemize}
\tightlist
\item
  State the replicator equation and explain the biological and
  game-theoretic intuition behind each term.
\item
  Use \texttt{deSolve} to numerically integrate the replicator equation
  for arbitrary payoff matrices.
\item
  Visualize replicator dynamics on the two-dimensional simplex for
  three-strategy games.
\item
  Analyze convergence to equilibrium in Hawk-Dove and cycling behavior
  in Rock-Paper-Scissors.
\end{itemize}

\subsection{Motivation}\label{motivation}

In @ref(sec-axelrod-tournament) and @ref(sec-spatial-pd), we studied
discrete agents making discrete choices. But many evolutionary and
economic processes are better described in continuous time: strategy
frequencies shift gradually as more successful types reproduce faster.
How do we model this?

The \textbf{replicator equation}, introduced by Taylor and Jonker
{[}-@taylor1978{]} and named by Schuster and Sigmund
{[}-@schuster1983{]}, provides the answer. It is the fundamental
dynamical system of evolutionary game theory, describing how the
frequency of each strategy in a large population changes over time as a
function of its fitness relative to the population average.

The replicator equation connects game theory to biology (it arises from
natural selection), economics (it describes learning dynamics in large
populations), and dynamical systems (it produces fixed points, limit
cycles, and chaos). In this chapter, we solve it numerically in R for
two classic games and visualize the resulting dynamics.

\subsection{Theory}\label{theory}

\subsubsection{The replicator equation}\label{the-replicator-equation}

Consider a large population where individuals play one of \(n\)
strategies. Let \(x_i\) be the fraction of the population playing
strategy \(i\), so \(\mathbf{x} = (x_1, \ldots, x_n)\) lies on the
\textbf{simplex}
\(\Delta^{n-1} = \{ \mathbf{x} \geq 0 : \sum_i x_i = 1 \}\).

The \textbf{fitness} of strategy \(i\) against the current population
is: \begin{equation}
f_i(\mathbf{x}) = \sum_j a_{ij} x_j = (A\mathbf{x})_i
$$
where $A = (a_{ij})$ is the payoff matrix. The **average fitness** is:
$$
\bar{f}(\mathbf{x}) = \sum_i x_i f_i(\mathbf{x}) = \mathbf{x}^\top A \mathbf{x}
$$

The **replicator equation** is:
\begin{equation}
\dot{x}_i = x_i \bigl( f_i(\mathbf{x}) - \bar{f}(\mathbf{x}) \bigr), \quad i = 1, \ldots, n
(\#eq:replicator)
\end{equation}

The interpretation is elegant: a strategy grows when its fitness exceeds
the population average, and shrinks when it falls below. The rate of
change is proportional to both the current frequency \(x_i\) (rare
strategies change slowly in absolute terms) and the fitness advantage
\(f_i - \bar{f}\).

\subsubsection{Key properties}\label{key-properties}

\begin{enumerate}
\def\labelenumi{\arabic{enumi}.}
\tightlist
\item
  \textbf{Simplex invariance}: If \(\mathbf{x}(0) \in \Delta^{n-1}\),
  then \(\mathbf{x}(t) \in \Delta^{n-1}\) for all \(t \geq 0\).
  Frequencies always sum to one.
\item
  \textbf{Rest points}: Every vertex of the simplex (a pure-strategy
  population) is a rest point. Interior rest points correspond to mixed
  Nash equilibria of the underlying game.
\item
  \textbf{Folk theorem}: In two-strategy games, every ESS is an
  asymptotically stable rest point of the replicator dynamics, and every
  stable rest point is a Nash equilibrium.
\end{enumerate}

\subsubsection{Two example games}\label{two-example-games}

\textbf{Rock-Paper-Scissors (RPS).} The payoff matrix is:

\[
A_{\text{RPS}} = \begin{pmatrix} 0 & -1 & 1 \\ 1 & 0 & -1 \\ -1 & 1 & 0 \end{pmatrix}
\]

The unique interior equilibrium \(\mathbf{x}^* = (1/3, 1/3, 1/3)\) is a
center: trajectories cycle around it forever without converging. This is
the classic example of non-convergent replicator dynamics.

\textbf{Hawk-Dove.} With benefit \(V = 2\) and cost \(C = 4\):

\[
A_{\text{HD}} = \begin{pmatrix} (V - C)/2 & V \\ 0 & V/2 \end{pmatrix} = \begin{pmatrix} -1 & 2 \\ 0 & 1 \end{pmatrix}
\]

The interior equilibrium is at \(x_H^* = V/C = 0.5\). This is an ESS and
the replicator dynamics converge to it from any interior initial
condition.

\subsection{Implementation in R}\label{sec-replicator-implementation}

\subsubsection{\texorpdfstring{Solving with
\texttt{run\_replicator()}}{Solving with run\_replicator()}}\label{solving-with-run_replicator}

The helper function \texttt{run\_replicator()} (defined in
\texttt{R/replicator.R}) wraps \texttt{deSolve::ode()} to integrate the
replicator equation. It takes a payoff matrix, initial frequencies, and
a time grid, returning a data frame of strategy frequencies over time.

\begin{Shaded}
\begin{Highlighting}[]
\NormalTok{A\_rps }\OtherTok{\textless{}{-}} \FunctionTok{matrix}\NormalTok{(}\FunctionTok{c}\NormalTok{( }\DecValTok{0}\NormalTok{, }\SpecialCharTok{{-}}\DecValTok{1}\NormalTok{,  }\DecValTok{1}\NormalTok{,}
                    \DecValTok{1}\NormalTok{,  }\DecValTok{0}\NormalTok{, }\SpecialCharTok{{-}}\DecValTok{1}\NormalTok{,}
                   \SpecialCharTok{{-}}\DecValTok{1}\NormalTok{,  }\DecValTok{1}\NormalTok{,  }\DecValTok{0}\NormalTok{), }\AttributeTok{nrow =} \DecValTok{3}\NormalTok{, }\AttributeTok{byrow =} \ConstantTok{TRUE}\NormalTok{)}

\NormalTok{A\_hd }\OtherTok{\textless{}{-}} \FunctionTok{matrix}\NormalTok{(}\FunctionTok{c}\NormalTok{(}\SpecialCharTok{{-}}\DecValTok{1}\NormalTok{, }\DecValTok{2}\NormalTok{,}
                   \DecValTok{0}\NormalTok{, }\DecValTok{1}\NormalTok{), }\AttributeTok{nrow =} \DecValTok{2}\NormalTok{, }\AttributeTok{byrow =} \ConstantTok{TRUE}\NormalTok{)}
\end{Highlighting}
\end{Shaded}

\subsubsection{Simplex coordinate
transform}\label{simplex-coordinate-transform}

To plot three-strategy dynamics on a two-dimensional ternary diagram, we
project from the simplex to Cartesian coordinates using the standard
transform:

\begin{Shaded}
\begin{Highlighting}[]
\NormalTok{simplex\_to\_cartesian }\OtherTok{\textless{}{-}} \ControlFlowTok{function}\NormalTok{(x1, x2, x3) \{}
  \FunctionTok{tibble}\NormalTok{(}
    \AttributeTok{sx =} \FloatTok{0.5} \SpecialCharTok{*}\NormalTok{ (}\DecValTok{2} \SpecialCharTok{*}\NormalTok{ x2 }\SpecialCharTok{+}\NormalTok{ x3) }\SpecialCharTok{/}\NormalTok{ (x1 }\SpecialCharTok{+}\NormalTok{ x2 }\SpecialCharTok{+}\NormalTok{ x3),}
    \AttributeTok{sy =}\NormalTok{ (}\FunctionTok{sqrt}\NormalTok{(}\DecValTok{3}\NormalTok{) }\SpecialCharTok{/} \DecValTok{2}\NormalTok{) }\SpecialCharTok{*}\NormalTok{ x3 }\SpecialCharTok{/}\NormalTok{ (x1 }\SpecialCharTok{+}\NormalTok{ x2 }\SpecialCharTok{+}\NormalTok{ x3)}
\NormalTok{  )}
\NormalTok{\}}
\end{Highlighting}
\end{Shaded}

\subsubsection{Running RPS dynamics from multiple initial
conditions}\label{running-rps-dynamics-from-multiple-initial-conditions}

\begin{Shaded}
\begin{Highlighting}[]
\FunctionTok{set.seed}\NormalTok{(}\DecValTok{42}\NormalTok{)}
\NormalTok{n\_traj }\OtherTok{\textless{}{-}} \DecValTok{6}
\NormalTok{initial\_conditions }\OtherTok{\textless{}{-}} \FunctionTok{map}\NormalTok{(}\FunctionTok{seq\_len}\NormalTok{(n\_traj), }\ControlFlowTok{function}\NormalTok{(i) \{}
\NormalTok{  raw }\OtherTok{\textless{}{-}} \FunctionTok{runif}\NormalTok{(}\DecValTok{3}\NormalTok{)}
\NormalTok{  raw }\SpecialCharTok{/} \FunctionTok{sum}\NormalTok{(raw)}
\NormalTok{\})}

\NormalTok{rps\_trajectories }\OtherTok{\textless{}{-}} \FunctionTok{map\_dfr}\NormalTok{(}\FunctionTok{seq\_along}\NormalTok{(initial\_conditions), }\ControlFlowTok{function}\NormalTok{(k) \{}
\NormalTok{  x0 }\OtherTok{\textless{}{-}}\NormalTok{ initial\_conditions[[k]]}
\NormalTok{  out }\OtherTok{\textless{}{-}} \FunctionTok{run\_replicator}\NormalTok{(A\_rps, x0, }\AttributeTok{times =} \FunctionTok{seq}\NormalTok{(}\DecValTok{0}\NormalTok{, }\DecValTok{40}\NormalTok{, }\AttributeTok{by =} \FloatTok{0.05}\NormalTok{))}
\NormalTok{  coords }\OtherTok{\textless{}{-}} \FunctionTok{simplex\_to\_cartesian}\NormalTok{(out[, }\DecValTok{2}\NormalTok{], out[, }\DecValTok{3}\NormalTok{], out[, }\DecValTok{4}\NormalTok{])}
\NormalTok{  coords}\SpecialCharTok{$}\NormalTok{trajectory }\OtherTok{\textless{}{-}} \FunctionTok{factor}\NormalTok{(k)}
\NormalTok{  coords}\SpecialCharTok{$}\NormalTok{time }\OtherTok{\textless{}{-}}\NormalTok{ out}\SpecialCharTok{$}\NormalTok{time}
\NormalTok{  coords}
\NormalTok{\})}
\end{Highlighting}
\end{Shaded}

\subsubsection{Running Hawk-Dove
dynamics}\label{running-hawk-dove-dynamics}

\begin{Shaded}
\begin{Highlighting}[]
\NormalTok{hd\_initials }\OtherTok{\textless{}{-}} \FunctionTok{c}\NormalTok{(}\FloatTok{0.1}\NormalTok{, }\FloatTok{0.3}\NormalTok{, }\FloatTok{0.5}\NormalTok{, }\FloatTok{0.7}\NormalTok{, }\FloatTok{0.9}\NormalTok{)}

\NormalTok{hd\_trajectories }\OtherTok{\textless{}{-}} \FunctionTok{map\_dfr}\NormalTok{(hd\_initials, }\ControlFlowTok{function}\NormalTok{(p) \{}
\NormalTok{  x0 }\OtherTok{\textless{}{-}} \FunctionTok{c}\NormalTok{(p, }\DecValTok{1} \SpecialCharTok{{-}}\NormalTok{ p)}
\NormalTok{  out }\OtherTok{\textless{}{-}} \FunctionTok{run\_replicator}\NormalTok{(A\_hd, x0, }\AttributeTok{times =} \FunctionTok{seq}\NormalTok{(}\DecValTok{0}\NormalTok{, }\DecValTok{20}\NormalTok{, }\AttributeTok{by =} \FloatTok{0.05}\NormalTok{))}
\NormalTok{  df }\OtherTok{\textless{}{-}} \FunctionTok{as\_tibble}\NormalTok{(out)}
  \FunctionTok{names}\NormalTok{(df) }\OtherTok{\textless{}{-}} \FunctionTok{c}\NormalTok{(}\StringTok{"time"}\NormalTok{, }\StringTok{"Hawk"}\NormalTok{, }\StringTok{"Dove"}\NormalTok{)}
\NormalTok{  df}\SpecialCharTok{$}\NormalTok{initial }\OtherTok{\textless{}{-}} \FunctionTok{paste0}\NormalTok{(}\StringTok{"x0 = "}\NormalTok{, p)}
\NormalTok{  df}
\NormalTok{\})}
\end{Highlighting}
\end{Shaded}

\subsection{Worked example}\label{worked-example}

We trace the replicator dynamics for specific initial conditions in both
games.

\subsubsection{Rock-Paper-Scissors}\label{rock-paper-scissors}

Starting from \(\mathbf{x}_0 = (0.6, 0.3, 0.1)\) --- a population
dominated by Rock:

\begin{Shaded}
\begin{Highlighting}[]
\NormalTok{x0\_rps }\OtherTok{\textless{}{-}} \FunctionTok{c}\NormalTok{(}\FloatTok{0.6}\NormalTok{, }\FloatTok{0.3}\NormalTok{, }\FloatTok{0.1}\NormalTok{)}
\NormalTok{out\_rps }\OtherTok{\textless{}{-}} \FunctionTok{run\_replicator}\NormalTok{(A\_rps, x0\_rps, }\AttributeTok{times =} \FunctionTok{seq}\NormalTok{(}\DecValTok{0}\NormalTok{, }\DecValTok{40}\NormalTok{, }\AttributeTok{by =} \FloatTok{0.1}\NormalTok{))}
\FunctionTok{names}\NormalTok{(out\_rps) }\OtherTok{\textless{}{-}} \FunctionTok{c}\NormalTok{(}\StringTok{"time"}\NormalTok{, }\StringTok{"Rock"}\NormalTok{, }\StringTok{"Paper"}\NormalTok{, }\StringTok{"Scissors"}\NormalTok{)}

\FunctionTok{cat}\NormalTok{(}\StringTok{"RPS dynamics from x0 = (0.6, 0.3, 0.1):}\SpecialCharTok{\textbackslash{}n\textbackslash{}n}\StringTok{"}\NormalTok{)}
\end{Highlighting}
\end{Shaded}

\begin{verbatim}
#> RPS dynamics from x0 = (0.6, 0.3, 0.1):
\end{verbatim}

\begin{Shaded}
\begin{Highlighting}[]
\FunctionTok{cat}\NormalTok{(}\FunctionTok{sprintf}\NormalTok{(}\StringTok{"  t = \%2.0f:  Rock = \%.3f  Paper = \%.3f  Scissors = \%.3f}\SpecialCharTok{\textbackslash{}n}\StringTok{"}\NormalTok{,}
            \DecValTok{0}\NormalTok{, x0\_rps[}\DecValTok{1}\NormalTok{], x0\_rps[}\DecValTok{2}\NormalTok{], x0\_rps[}\DecValTok{3}\NormalTok{]))}
\end{Highlighting}
\end{Shaded}

\begin{verbatim}
#>   t =  0:  Rock = 0.600  Paper = 0.300  Scissors = 0.100
\end{verbatim}

\begin{Shaded}
\begin{Highlighting}[]
\ControlFlowTok{for}\NormalTok{ (t\_val }\ControlFlowTok{in} \FunctionTok{c}\NormalTok{(}\DecValTok{5}\NormalTok{, }\DecValTok{10}\NormalTok{, }\DecValTok{20}\NormalTok{, }\DecValTok{40}\NormalTok{)) \{}
\NormalTok{  row }\OtherTok{\textless{}{-}}\NormalTok{ out\_rps[out\_rps}\SpecialCharTok{$}\NormalTok{time }\SpecialCharTok{==}\NormalTok{ t\_val, ]}
  \FunctionTok{cat}\NormalTok{(}\FunctionTok{sprintf}\NormalTok{(}\StringTok{"  t = \%2.0f:  Rock = \%.3f  Paper = \%.3f  Scissors = \%.3f}\SpecialCharTok{\textbackslash{}n}\StringTok{"}\NormalTok{,}
\NormalTok{              t\_val, row}\SpecialCharTok{$}\NormalTok{Rock, row}\SpecialCharTok{$}\NormalTok{Paper, row}\SpecialCharTok{$}\NormalTok{Scissors))}
\NormalTok{\}}
\end{Highlighting}
\end{Shaded}

\begin{verbatim}
#>   t =  5:  Rock = 0.087  Paper = 0.488  Scissors = 0.426
#>   t = 10:  Rock = 0.551  Paper = 0.091  Scissors = 0.358
#>   t = 20:  Rock = 0.164  Paper = 0.163  Scissors = 0.673
#>   t = 40:  Rock = 0.298  Paper = 0.602  Scissors = 0.100
\end{verbatim}

The frequencies cycle endlessly around the interior equilibrium
\((1/3, 1/3, 1/3)\). Rock starts high, attracting Paper players who can
exploit it; Paper then rises, attracting Scissors; Scissors rises,
attracting Rock --- and the cycle repeats. The orbit never converges
because the interior equilibrium is a center (neutrally stable), not an
attractor.

\begin{Shaded}
\begin{Highlighting}[]
\NormalTok{triangle }\OtherTok{\textless{}{-}} \FunctionTok{tibble}\NormalTok{(}
  \AttributeTok{sx =} \FunctionTok{c}\NormalTok{(}\DecValTok{0}\NormalTok{, }\DecValTok{1}\NormalTok{, }\FloatTok{0.5}\NormalTok{, }\DecValTok{0}\NormalTok{),}
  \AttributeTok{sy =} \FunctionTok{c}\NormalTok{(}\DecValTok{0}\NormalTok{, }\DecValTok{0}\NormalTok{, }\FunctionTok{sqrt}\NormalTok{(}\DecValTok{3}\NormalTok{) }\SpecialCharTok{/} \DecValTok{2}\NormalTok{, }\DecValTok{0}\NormalTok{)}
\NormalTok{)}

\NormalTok{center }\OtherTok{\textless{}{-}} \FunctionTok{simplex\_to\_cartesian}\NormalTok{(}\DecValTok{1}\SpecialCharTok{/}\DecValTok{3}\NormalTok{, }\DecValTok{1}\SpecialCharTok{/}\DecValTok{3}\NormalTok{, }\DecValTok{1}\SpecialCharTok{/}\DecValTok{3}\NormalTok{)}

\NormalTok{p\_rps }\OtherTok{\textless{}{-}} \FunctionTok{ggplot}\NormalTok{() }\SpecialCharTok{+}
  \FunctionTok{geom\_path}\NormalTok{(}\AttributeTok{data =}\NormalTok{ triangle, }\FunctionTok{aes}\NormalTok{(}\AttributeTok{x =}\NormalTok{ sx, }\AttributeTok{y =}\NormalTok{ sy),}
            \AttributeTok{colour =} \StringTok{"grey40"}\NormalTok{, }\AttributeTok{linewidth =} \FloatTok{0.5}\NormalTok{) }\SpecialCharTok{+}
  \FunctionTok{geom\_path}\NormalTok{(}\AttributeTok{data =}\NormalTok{ rps\_trajectories,}
            \FunctionTok{aes}\NormalTok{(}\AttributeTok{x =}\NormalTok{ sx, }\AttributeTok{y =}\NormalTok{ sy, }\AttributeTok{colour =}\NormalTok{ trajectory),}
            \AttributeTok{linewidth =} \FloatTok{0.5}\NormalTok{, }\AttributeTok{alpha =} \FloatTok{0.8}\NormalTok{) }\SpecialCharTok{+}
  \FunctionTok{geom\_point}\NormalTok{(}\AttributeTok{data =}\NormalTok{ center, }\FunctionTok{aes}\NormalTok{(}\AttributeTok{x =}\NormalTok{ sx, }\AttributeTok{y =}\NormalTok{ sy),}
             \AttributeTok{shape =} \DecValTok{3}\NormalTok{, }\AttributeTok{size =} \DecValTok{3}\NormalTok{, }\AttributeTok{stroke =} \FloatTok{1.2}\NormalTok{, }\AttributeTok{colour =}\NormalTok{ okabe\_ito[}\DecValTok{8}\NormalTok{]) }\SpecialCharTok{+}
  \FunctionTok{annotate}\NormalTok{(}\StringTok{"text"}\NormalTok{, }\AttributeTok{x =} \SpecialCharTok{{-}}\FloatTok{0.05}\NormalTok{, }\AttributeTok{y =} \SpecialCharTok{{-}}\FloatTok{0.04}\NormalTok{, }\AttributeTok{label =} \StringTok{"Rock"}\NormalTok{,}
           \AttributeTok{size =} \FloatTok{3.5}\NormalTok{, }\AttributeTok{fontface =} \StringTok{"bold"}\NormalTok{) }\SpecialCharTok{+}
  \FunctionTok{annotate}\NormalTok{(}\StringTok{"text"}\NormalTok{, }\AttributeTok{x =} \FloatTok{1.05}\NormalTok{, }\AttributeTok{y =} \SpecialCharTok{{-}}\FloatTok{0.04}\NormalTok{, }\AttributeTok{label =} \StringTok{"Paper"}\NormalTok{,}
           \AttributeTok{size =} \FloatTok{3.5}\NormalTok{, }\AttributeTok{fontface =} \StringTok{"bold"}\NormalTok{) }\SpecialCharTok{+}
  \FunctionTok{annotate}\NormalTok{(}\StringTok{"text"}\NormalTok{, }\AttributeTok{x =} \FloatTok{0.5}\NormalTok{, }\AttributeTok{y =} \FunctionTok{sqrt}\NormalTok{(}\DecValTok{3}\NormalTok{)}\SpecialCharTok{/}\DecValTok{2} \SpecialCharTok{+} \FloatTok{0.04}\NormalTok{, }\AttributeTok{label =} \StringTok{"Scissors"}\NormalTok{,}
           \AttributeTok{size =} \FloatTok{3.5}\NormalTok{, }\AttributeTok{fontface =} \StringTok{"bold"}\NormalTok{) }\SpecialCharTok{+}
  \FunctionTok{scale\_colour\_manual}\NormalTok{(}\AttributeTok{values =}\NormalTok{ okabe\_ito[}\DecValTok{1}\SpecialCharTok{:}\DecValTok{6}\NormalTok{]) }\SpecialCharTok{+}
  \FunctionTok{coord\_equal}\NormalTok{(}\AttributeTok{xlim =} \FunctionTok{c}\NormalTok{(}\SpecialCharTok{{-}}\FloatTok{0.1}\NormalTok{, }\FloatTok{1.1}\NormalTok{), }\AttributeTok{ylim =} \FunctionTok{c}\NormalTok{(}\SpecialCharTok{{-}}\FloatTok{0.1}\NormalTok{, }\FloatTok{1.0}\NormalTok{)) }\SpecialCharTok{+}
  \FunctionTok{theme\_publication}\NormalTok{() }\SpecialCharTok{+}
  \FunctionTok{theme}\NormalTok{(}\AttributeTok{axis.text =} \FunctionTok{element\_blank}\NormalTok{(),}
        \AttributeTok{axis.ticks =} \FunctionTok{element\_blank}\NormalTok{(),}
        \AttributeTok{panel.grid.major =} \FunctionTok{element\_blank}\NormalTok{(),}
        \AttributeTok{legend.position =} \StringTok{"none"}\NormalTok{) }\SpecialCharTok{+}
  \FunctionTok{labs}\NormalTok{(}\AttributeTok{title =} \StringTok{"RPS Replicator Dynamics on the Simplex"}\NormalTok{,}
       \AttributeTok{x =} \ConstantTok{NULL}\NormalTok{, }\AttributeTok{y =} \ConstantTok{NULL}\NormalTok{)}

\NormalTok{p\_rps}
\end{Highlighting}
\end{Shaded}

\begin{figure}

{\centering \includegraphics[width=0.8\linewidth]{/home/runner/work/strategy-in-r/strategy-in-r/docs/pdf-chapters/19-replicator-dynamics_files/figure-latex/rps-simplex-1} 

}

\caption{Replicator dynamics for Rock-Paper-Scissors on the two-dimensional simplex. Six trajectories from random initial conditions cycle around the interior equilibrium at (1/3, 1/3, 1/3). The orbits are closed curves — the system never converges, reflecting the non-transitive dominance structure of the game.}\label{fig:rps-simplex}
\end{figure}

\begin{Shaded}
\begin{Highlighting}[]
\FunctionTok{save\_pub\_fig}\NormalTok{(p\_rps, }\StringTok{"replicator{-}rps{-}simplex"}\NormalTok{, }\AttributeTok{width =} \DecValTok{7}\NormalTok{, }\AttributeTok{height =} \DecValTok{6}\NormalTok{)}
\end{Highlighting}
\end{Shaded}

\subsubsection{Hawk-Dove}\label{hawk-dove}

Starting from \(x_H = 0.1\) (a population with few Hawks):

\begin{Shaded}
\begin{Highlighting}[]
\NormalTok{x0\_hd }\OtherTok{\textless{}{-}} \FunctionTok{c}\NormalTok{(}\FloatTok{0.1}\NormalTok{, }\FloatTok{0.9}\NormalTok{)}
\NormalTok{out\_hd }\OtherTok{\textless{}{-}} \FunctionTok{run\_replicator}\NormalTok{(A\_hd, x0\_hd, }\AttributeTok{times =} \FunctionTok{seq}\NormalTok{(}\DecValTok{0}\NormalTok{, }\DecValTok{20}\NormalTok{, }\AttributeTok{by =} \FloatTok{0.1}\NormalTok{))}
\FunctionTok{names}\NormalTok{(out\_hd) }\OtherTok{\textless{}{-}} \FunctionTok{c}\NormalTok{(}\StringTok{"time"}\NormalTok{, }\StringTok{"Hawk"}\NormalTok{, }\StringTok{"Dove"}\NormalTok{)}

\FunctionTok{cat}\NormalTok{(}\StringTok{"Hawk{-}Dove dynamics from x0 = (0.1, 0.9):}\SpecialCharTok{\textbackslash{}n\textbackslash{}n}\StringTok{"}\NormalTok{)}
\end{Highlighting}
\end{Shaded}

\begin{verbatim}
#> Hawk-Dove dynamics from x0 = (0.1, 0.9):
\end{verbatim}

\begin{Shaded}
\begin{Highlighting}[]
\ControlFlowTok{for}\NormalTok{ (t\_val }\ControlFlowTok{in} \FunctionTok{c}\NormalTok{(}\DecValTok{0}\NormalTok{, }\DecValTok{2}\NormalTok{, }\DecValTok{5}\NormalTok{, }\DecValTok{10}\NormalTok{, }\DecValTok{20}\NormalTok{)) \{}
\NormalTok{  row }\OtherTok{\textless{}{-}}\NormalTok{ out\_hd[out\_hd}\SpecialCharTok{$}\NormalTok{time }\SpecialCharTok{==}\NormalTok{ t\_val, ]}
  \FunctionTok{cat}\NormalTok{(}\FunctionTok{sprintf}\NormalTok{(}\StringTok{"  t = \%2.0f:  Hawk = \%.3f  Dove = \%.3f}\SpecialCharTok{\textbackslash{}n}\StringTok{"}\NormalTok{,}
\NormalTok{              t\_val, row}\SpecialCharTok{$}\NormalTok{Hawk, row}\SpecialCharTok{$}\NormalTok{Dove))}
\NormalTok{\}}
\end{Highlighting}
\end{Shaded}

\begin{verbatim}
#>   t =  0:  Hawk = 0.100  Dove = 0.900
#>   t =  2:  Hawk = 0.280  Dove = 0.720
#>   t =  5:  Hawk = 0.446  Dove = 0.554
#>   t = 10:  Hawk = 0.496  Dove = 0.504
#>   t = 20:  Hawk = 0.500  Dove = 0.500
\end{verbatim}

Hawks are rare but earn a high payoff by exploiting the predominantly
Dove population. The Hawk frequency rises rapidly at first, then
decelerates as more Hawk-Hawk encounters occur (which are costly). The
system converges smoothly to the ESS at \(x_H^* = V/C = 0.5\).

\begin{Shaded}
\begin{Highlighting}[]
\NormalTok{p\_hd }\OtherTok{\textless{}{-}} \FunctionTok{ggplot}\NormalTok{(hd\_trajectories, }\FunctionTok{aes}\NormalTok{(}\AttributeTok{x =}\NormalTok{ time, }\AttributeTok{y =}\NormalTok{ Hawk, }\AttributeTok{colour =}\NormalTok{ initial)) }\SpecialCharTok{+}
  \FunctionTok{geom\_line}\NormalTok{(}\AttributeTok{linewidth =} \FloatTok{0.8}\NormalTok{) }\SpecialCharTok{+}
  \FunctionTok{geom\_hline}\NormalTok{(}\AttributeTok{yintercept =} \FloatTok{0.5}\NormalTok{, }\AttributeTok{linetype =} \StringTok{"dashed"}\NormalTok{, }\AttributeTok{colour =}\NormalTok{ okabe\_ito[}\DecValTok{8}\NormalTok{]) }\SpecialCharTok{+}
  \FunctionTok{annotate}\NormalTok{(}\StringTok{"text"}\NormalTok{, }\AttributeTok{x =} \DecValTok{18}\NormalTok{, }\AttributeTok{y =} \FloatTok{0.53}\NormalTok{, }\AttributeTok{label =} \StringTok{"ESS: x[H] == 0.5"}\NormalTok{,}
           \AttributeTok{parse =} \ConstantTok{TRUE}\NormalTok{, }\AttributeTok{size =} \FloatTok{3.5}\NormalTok{, }\AttributeTok{colour =}\NormalTok{ okabe\_ito[}\DecValTok{8}\NormalTok{]) }\SpecialCharTok{+}
  \FunctionTok{scale\_colour\_manual}\NormalTok{(}\AttributeTok{values =}\NormalTok{ okabe\_ito[}\DecValTok{1}\SpecialCharTok{:}\DecValTok{5}\NormalTok{], }\AttributeTok{name =} \StringTok{"Initial Hawk freq."}\NormalTok{) }\SpecialCharTok{+}
  \FunctionTok{scale\_y\_continuous}\NormalTok{(}\AttributeTok{limits =} \FunctionTok{c}\NormalTok{(}\DecValTok{0}\NormalTok{, }\DecValTok{1}\NormalTok{), }\AttributeTok{labels =} \FunctionTok{label\_percent}\NormalTok{()) }\SpecialCharTok{+}
  \FunctionTok{theme\_publication}\NormalTok{() }\SpecialCharTok{+}
  \FunctionTok{labs}\NormalTok{(}\AttributeTok{title =} \StringTok{"Hawk{-}Dove Replicator Dynamics"}\NormalTok{,}
       \AttributeTok{x =} \StringTok{"Time"}\NormalTok{, }\AttributeTok{y =} \StringTok{"Hawk frequency"}\NormalTok{)}

\NormalTok{p\_hd}
\end{Highlighting}
\end{Shaded}

\begin{figure}

{\centering \includegraphics[width=0.8\linewidth]{/home/runner/work/strategy-in-r/strategy-in-r/docs/pdf-chapters/19-replicator-dynamics_files/figure-latex/hd-dynamics-1} 

}

\caption{Hawk frequency over time for five different initial conditions. All trajectories converge to the ESS at x_H = 0.5 (dashed line), regardless of starting point. Initial conditions below the ESS see Hawk increase; those above see it decrease.}\label{fig:hd-dynamics}
\end{figure}

\begin{Shaded}
\begin{Highlighting}[]
\FunctionTok{save\_pub\_fig}\NormalTok{(p\_hd, }\StringTok{"replicator{-}hd{-}dynamics"}\NormalTok{, }\AttributeTok{width =} \DecValTok{7}\NormalTok{, }\AttributeTok{height =} \DecValTok{5}\NormalTok{)}
\end{Highlighting}
\end{Shaded}

The Hawk-Dove dynamics illustrate a core result in evolutionary game
theory: \textbf{the ESS is a global attractor of the replicator
dynamics}. No matter how the population starts (as long as both
strategies are present), it converges to the mixed equilibrium where
Hawks and Doves coexist. This is in stark contrast to RPS, where the
dynamics cycle forever.

\subsubsection{Comparing the two games}\label{comparing-the-two-games}

The difference between the two games reflects a deep structural
distinction:

\begin{itemize}
\tightlist
\item
  \textbf{Hawk-Dove} has a symmetric \(2 \times 2\) payoff matrix with
  an interior ESS. The replicator dynamics have a Lyapunov function (a
  quantity that strictly decreases along trajectories), guaranteeing
  convergence.
\item
  \textbf{Rock-Paper-Scissors} has a skew-symmetric payoff matrix. The
  quantity \(H(\mathbf{x}) = x_1 x_2 x_3\) is conserved along
  trajectories (it is a constant of motion), producing closed orbits
  rather than convergence.
\end{itemize}

This shows that the replicator equation's behavior depends qualitatively
on the game's payoff structure --- a theme we will revisit when studying
evolutionarily stable strategies in @ref(sec-ess).

\subsection{Extensions}\label{extensions}

\begin{itemize}
\tightlist
\item
  \textbf{Asymmetric replicator dynamics}: When players occupy different
  roles (e.g., owner vs.~intruder), the dynamics involve two coupled
  replicator equations, one for each population. This leads to
  higher-dimensional phase spaces.
\item
  \textbf{Mutation and selection}: Adding a small mutation rate \(\mu\)
  transforms the replicator equation into the \textbf{replicator-mutator
  equation}: \(\dot{x}_i = \sum_j x_j f_j Q_{ji} - x_i \bar{f}\), where
  \(Q\) is a mutation matrix. This can turn centers into spirals.
\item
  \textbf{Discrete-time replicator}: The discrete-time version
  \(x_i' = x_i f_i / \bar{f}\) has different stability properties and
  can exhibit chaos even in low dimensions.
\item
  \textbf{Connection to spatial models}: The well-mixed replicator
  dynamics serve as a baseline for the spatial models of
  @ref(sec-spatial-pd). When spatial structure matters, local
  frequency-dependent selection can sustain strategies that the
  replicator equation predicts will go extinct.
\item
  For a comprehensive treatment, see @hofbauer1998. For the connection
  between replicator dynamics and learning, see @fudenberg1998.
\end{itemize}

\subsection*{Exercises}\label{exercises}
\addcontentsline{toc}{subsection}{Exercises}

\begin{enumerate}
\def\labelenumi{\arabic{enumi}.}
\item
  \textbf{Three-strategy Hawk-Dove-Retaliator.} Define a \(3 \times 3\)
  payoff matrix for a game with Hawks, Doves, and Retaliators
  (Retaliators play Dove against Doves and Hawk against Hawks). With
  \(V = 2\) and \(C = 4\), the matrix is:

  \[A = \begin{pmatrix} -1 & 2 & -1 \\ 0 & 1 & 1 \\ -1 & 1 & 0.5 \end{pmatrix}\]

  Use \texttt{run\_replicator()} to solve the dynamics from three
  different initial conditions. Plot the trajectories on the simplex.
  Does the system converge? If so, to what equilibrium?
\item
  \textbf{Sensitivity to initial conditions in RPS.} Run the RPS
  dynamics from \(\mathbf{x}_0 = (0.5, 0.3, 0.2)\) and from
  \(\mathbf{x}_0 = (0.501, 0.3, 0.199)\). Plot both trajectories on the
  simplex. How different are the orbits? Relate your finding to the
  conservation of \(H(\mathbf{x}) = x_1 x_2 x_3\).
\item
  \textbf{Payoff perturbation.} Modify the RPS payoff matrix to be
  slightly asymmetric: change the (1,2) entry from \(-1\) to \(-1.1\).
  Run the dynamics and plot the trajectory. Does the system still cycle,
  or does it now converge to a vertex? Explain why breaking the
  skew-symmetry changes the qualitative dynamics.
\end{enumerate}

Solutions appear in @ref(sec-solutions).

\end{document}
